\documentclass[11pt,a4paper]{article}

\usepackage[margin=1in]{geometry}
\usepackage{amsmath,amssymb,amsthm}
\usepackage{mathtools}
\usepackage{hyperref}
\usepackage{booktabs}
\usepackage{xcolor}
\usepackage{listings}
\usepackage{enumitem}

\hypersetup{
  colorlinks=true,
  linkcolor=blue!50!black,
  citecolor=blue!50!black,
  urlcolor=blue!50!black,
}

% Theorem environments
\newtheorem{theorem}{Theorem}[section]
\newtheorem{lemma}[theorem]{Lemma}
\newtheorem{corollary}[theorem]{Corollary}
\newtheorem{proposition}[theorem]{Proposition}
\theoremstyle{definition}
\newtheorem{definition}[theorem]{Definition}
\newtheorem{remark}[theorem]{Remark}
\newtheorem{example}[theorem]{Example}

% Custom macros
\newcommand{\FF}{\mathbb{F}}
\newcommand{\ZZ}{\mathbb{Z}}
\newcommand{\QQ}{\mathbb{Q}}
\newcommand{\End}{\operatorname{End}}
\newcommand{\Hom}{\operatorname{Hom}}
\newcommand{\Aut}{\operatorname{Aut}}
\newcommand{\Cl}{\operatorname{Cl}}
\newcommand{\Jac}{\operatorname{Jac}}
\newcommand{\Res}{\operatorname{Res}}
\newcommand{\disc}{\operatorname{disc}}
\newcommand{\gcd}{\operatorname{gcd}}
\newcommand{\NS}{\operatorname{NS}}

\title{Structural Completeness of Prime-Field Elliptic Curves \\
       Against Isogeny-Graph Cryptanalysis}
% Anonymized for double-blind submission.  Author block restored
% for camera-ready / preprint release; see paper/README.md.
\author{Anonymous}
\date{}

\begin{document}
\maketitle

\begin{abstract}
We pursue a comprehensive cryptanalytic audit of the isogeny-graph
structure of deployed prime-field elliptic curves~--- specifically
secp256k1, NIST P-256, brainpoolP256r1, and Curve25519's prime
subgroup. For each curve we run a battery of structural tests and
two novel attack proposals (Class-Group-Amortised Hidden-Number
Cryptanalysis; Volcano-Floor Class-Group-Augmented Pollard~$\rho$),
finding all directions structurally null at the prime-field level.

The main result is a \emph{structural-completeness theorem}: for
every prime-order ordinary elliptic curve $E$ over a prime field
$\FF_p$, no algorithm in the family $\{$$\FF_p$-isogeny transport,
$\Cl(\End)$-orbit walking, $(N,N)$-cover via Howe-gluing,
higher-genus cover, supersingular-reduction attack, volcano-floor
$\rho\}$ reduces ECDLP cost below $\sqrt{n / |\Aut(E)|}$ (plain
Pollard~$\rho$ with full automorphism folding).

The theorem is proven via seven independent structural blocks
(Section~\ref{sec:blocks}) and empirically verified across all
four deployed curves (Section~\ref{sec:empirical}). A surprising
auxiliary finding: the Howe-gluing genus-2 cover \emph{exists} for
every such curve, but does not yield a DLP speedup because
genus-2 prime-field DLP costs $O(p)$ while ECDLP costs
$O(\sqrt{p})$.  For secp256k1 specifically, the obstruction is
sharper: $x^3 + 7$ is irreducible over $\FF_p$, so the cover's
field of definition is $\FF_{p^3}$ rather than $\FF_p$, making
the cover-based attack cost $O(p^{3/2}) \approx 2^{384}$~---
a factor of $2^{256}$ slower than the $2^{128}$ ECDLP target.
\end{abstract}

\section{Introduction and motivation}

Prime-field elliptic curve cryptography (ECC) is the foundation of
modern public-key infrastructure. Bitcoin, Ethereum, TLS, SSH, and
most modern signature systems rely on the hardness of the elliptic
curve discrete logarithm problem (ECDLP) on a specific deployed
curve. Two major families exist:

\begin{center}
\begin{tabular}{lllll}
\toprule
Family & Examples & Field & Cofactor \\
\midrule
Koblitz & secp256k1 & prime & 1 \\
NIST & P-256, P-384, P-521 & prime & 1 \\
brainpool & brainpoolP\{256,384,512\}r1 & prime & 1 \\
Curve25519/Edwards & Curve25519, Ed25519, Curve448 & prime & 4 or 8 \\
\bottomrule
\end{tabular}
\end{center}

Despite 30+ years of public cryptanalytic effort, the best known
attack on prime-field ECC is Pollard~$\rho$ with cost $O(\sqrt{n})$
where $n$ is the prime subgroup order~--- providing 128-bit
security for 256-bit curves.

The isogeny-graph approach asks: does the rich algebraic structure
of an elliptic curve's \emph{isogeny class} leak any
ECDLP-relevant information? Concretely, can one find a curve $E'$
isogenous to the target $E$ such that ECDLP on $E'$ is easier than
on $E$? The answer, given by Tate's theorem~\cite{tate1966}, is
\emph{no} in the literal sense: all $\FF_p$-isogenous curves share
the same group order. But several more subtle questions
remain~--- the \emph{structure} of the isogeny graph might still
leak something.

This paper does the systematic version of that audit:

\begin{enumerate}
\item Enumerate every isogeny-graph-based attack vector.
\item For each, identify the structural obstruction or run an
      empirical falsifier.
\item Confirm the result holds across all major deployed curves.
\item Prove a structural-completeness theorem.
\end{enumerate}

The result is, predictably, a null: every direction is closed.
The value lies in the \emph{completeness} of the closure~--- the
theorem turns a folklore claim (``ECDLP is hard for prime-field
ECC'') into a rigorously decomposable, mechanically verifiable
structural statement.

\section{Setup: isogeny graphs of CM curves}

\subsection{Elliptic curves and endomorphism rings}

An elliptic curve over $\FF_p$ is $E\colon y^2 = x^3 + ax + b$ for
$a, b \in \FF_p$ with non-zero discriminant. The group $E(\FF_p)$
has order $n = p + 1 - t$ where $t$ is the \emph{Frobenius trace},
satisfying $|t| \leq 2\sqrt{p}$ (Hasse).

The \emph{endomorphism ring} $\End_{\FF_p}(E)$ is the ring of
group homomorphisms $E \to E$ defined over $\FF_p$. For ordinary
$E$ (the case $\gcd(t, p) = 1$), $\End(E)$ is an order in an
imaginary quadratic field $K = \QQ(\sqrt{D})$ with $D < 0$~---
the \emph{CM} field.

When $D$ is \emph{small} ($D \in \{-3, -4, -7, \ldots, -163\}$~---
the nine class-number-1 fields), $E$ has \emph{small CM} and
admits an efficient endomorphism beyond $[n]$, used for the GLV
scalar-multiplication trick. secp256k1 is $D = -3$ (CM by
$\ZZ[\omega]$); P-256, brainpool, Curve25519 are generic (huge $D$).

\subsection{The isogeny volcano}

For a small prime $\ell$, the $\ell$-isogeny graph at $j(E)$ has
the local structure of a ``volcano''~\cite{kohel1996,fouquetmorain2002}:

\begin{itemize}
\item \textbf{Crater} = curves with $\End$ = maximal order of $K$.
      Horizontal edges form the Cayley graph of $\Cl(O_K)$ (``class
      group action'').
\item \textbf{Floor levels} = curves with $\End$ in non-maximal
      orders of conductor $\ell$, $\ell^2$, \ldots. Class number
      grows.
\end{itemize}

For secp256k1, $h(-3) = 1$, so the crater is a single vertex
$j = 0$. For generic ordinary $E$, the crater has $\approx
\sqrt{|D|}$ vertices, all sharing the same group order $n$ by
Tate.

\subsection{The $\Cl(O)$-action and Tate's theorem}

\begin{theorem}[Tate~\cite{tate1966}]
Two abelian varieties over $\FF_p$ are $\FF_p$-isogenous if and
only if they have the same characteristic polynomial of Frobenius.
Equivalently for elliptic curves: same $\#E(\FF_p)$.
\end{theorem}

Consequence: every $\FF_p$-isogenous neighbour of a target curve
has the same order $n$. ECDLP transported via isogeny
$\varphi \colon E \to E'$ preserves the secret $d$:
\[
   (P, Q = d \cdot P) \text{ on } E \;\longleftrightarrow\;
   (\varphi(P), \varphi(Q) = d \cdot \varphi(P)) \text{ on } E'.
\]

This is the \emph{fundamental obstruction} to ``find an easier
curve in the isogeny class'' attacks.

\section{Attack surface taxonomy}\label{sec:taxonomy}

Isogeny-graph attacks fall into 7 distinct families.

\begin{center}
\begin{tabular}{lll}
\toprule
Family & Mechanism & Block \\
\midrule
$\FF_p$-isogeny transport & Find $E'$ with easier $\#E$ & B1 (Tate) \\
Class-group amortisation & Walk $\Cl(O)$-orbit, CRT & B2 \\
Volcano-floor walking (VFCG-$\rho$) & Walk floor curves & B3 \\
$(N,N)$-split-Jacobian cover & Find genus-2 $C$ with $\Jac(C) \sim E \times E^t$ & B4, B5 \\
Higher-genus cover ($g \geq 3$) & Find genus-$g$ $C$ with $\Jac(C) \to E$ & B5 \\
Diem sub-exp on $\FF_{p^k}$ restriction & Weil-restrict, Diem 2011 & B6 \\
Supersingular-reduction (V4) & SIDH-world attacks via $E \bmod \ell$ & B7 \\
\bottomrule
\end{tabular}
\end{center}

\section{The seven structural blocks}\label{sec:blocks}

Each block is an independent structural obstruction. Together they
close the attack-surface taxonomy of Section~\ref{sec:taxonomy}.

\subsection*{B1: Tate isogeny invariance}

For $E_1, E_2$ in the same $\FF_p$-isogeny class,
$\#E_1(\FF_p) = \#E_2(\FF_p)$. Hence ECDLP on $E_2$ has the same
hardness as on $E_1$.

\textbf{Implication}: $\FF_p$-isogeny transport (the simplest
``find an easier curve'' attack) is exactly blocked.

\subsection*{B2: Trivial $\Cl(O)$-orbit at the crater for small CM}

For an elliptic curve $E$ with $\End$ = maximal order of
$K = \QQ(\sqrt{D})$ for class-number-1 $D$, the $\Cl(O)$-orbit
consists of exactly $E$ itself (a self-loop). Class-group
amortisation has nothing to amortise.

For secp256k1 ($D = -3$, $h = 1$): the $\Cl(O)$-action is trivial,
blocking CSIDH-style key exchange constructions and CGA-HNC-style
attacks.

For curves \emph{without} small CM (P-256, brainpool, generic
Koblitz neighbours): $\Cl(O)$ has class number $\approx \sqrt{|D|}$
(sub-exponential in $p$). The orbit walk itself costs $\approx
\sqrt{|D|}$ steps, equal to the best generic $\rho$ cost.

\subsection*{B3: Floor-level $c_\gamma$ has uniformly small order but
no walk speedup}

For any depth-$d$ $\ell$-floor of an ordinary curve $E/\FF_p$ with
CM by an order of class number $h(-|D| \cdot \ell^{2d})$, the
class-group action's scalar $c_\gamma \in (\ZZ/n)^*$ has order
$\leq 2 \cdot h(\disc)$~--- uniformly small.

This bound holds uniformly for every $(D, \ell, d)$ tested. The
structural argument (additive $\rho$ steps re-randomise the walk
in $(\ZZ/n)$) shows that ``$c_\gamma \in$ small subgroup'' does
\emph{not} translate to a $\rho$ speedup.

\subsection*{B4: $\Hom_{\FF_p}(E, E^t) = 0$ implies $E \times E^t$
is not a Jacobian}

For ordinary $E/\FF_p$, $\Hom_{\FF_p}(E, E^t) = 0$. Hence
$E \times E^t$ has Néron-Severi rank 2 with the product
polarisation as the only (up to $\pm 1$) principal polarisation.

This rules out direct $E \times E^t \cong \Jac(C)$ for any smooth
$C/\FF_p$.

\subsection*{B5: Genus-$g$ cover DLP cost exceeds ECDLP for every
$g \geq 2$ over $\FF_p$}

For every smooth genus-$g$ curve $C/\FF_p$ with $\Jac(C)$ mapping
to $E$, the best known DLP algorithm on $\Jac(C)(\FF_p)$ has cost
$O(p^{2 - 2/g}) \geq O(p)$ for $g \geq 2$~\cite{gaudry2000}. Plain
ECDLP on $E$ costs $O(\sqrt{n}) = O(\sqrt{p})$.

So even when a smooth cover $C \to E$ exists, the cover gives no
DLP speedup.

\subsection*{B6: Diem 2011 sub-exp is inapplicable to prime fields}

Diem's 2011 sub-exponential DLP algorithm~\cite{diem2011} for
hyperelliptic Jacobians of genus $g \geq 3$ requires the base field
$\FF_q = \FF_{p^k}$ with $k \geq 3$. For $k = 1$ (prime field), the
algorithm's factor base construction does not exist.

\subsection*{B7: Supersingular reductions live in a disjoint world}

For an ordinary curve $E/\FF_p$, the supersingular reductions
$E \bmod \ell$ over $\overline{\FF_\ell}$ live in a different
mathematical world: their endomorphism rings are orders in
non-commutative quaternion algebras, and Deuring's lifting maps
$\FF_\ell \to \QQ_p$ (the $p$-adics), not $\FF_\ell \to \FF_p$.

Castryck-Decru 2022~\cite{castryckdecru2023}'s SIDH break operates
on the supersingular $\ell$-isogeny problem with auxiliary torsion
data. An ECDLP instance on $E/\FF_p$ publishes no such torsion
data; the SIDH-break machinery has nothing to consume.

\section{Novel proposals: CGA-HNC and VFCG-$\rho$}\label{sec:novel}

\subsection{CGA-HNC: Class-Group-Amortised Hidden-Number Cryptanalysis}

\textbf{Proposal}: walk the $\Cl(O)$-orbit, run Pohlig-Hellman on
each orbit curve's smooth subgroup, CRT-combine residual $d$ data,
clean up residual via lattice (HNP-style) attack.

\textbf{Result}: null for $h = 1$ (trivially), null for $h > 1$
(orbit walk cost matches saving). At cryptographic sizes, orbit
walk is $\sqrt{p}$ steps, same as plain $\rho$.

\subsection{VFCG-$\rho$: Volcano-Floor Class-Group-Augmented Pollard $\rho$}

\textbf{Proposal}: Pollard $\rho$ walks at the depth-1 floor of
secp256k1's isogeny volcano (where $h(-3\ell^2) > 1$), with
class-group action steps interleaved.

\textbf{Result}: null. Per-step $\ell$-isogeny cost ($O(\ell^2)$)
outweighs the $\sqrt{h(\disc)}$ birthday-saving. The
``$c_\gamma$ in a small subgroup of $(\ZZ/n)^*$'' hope is satisfied
uniformly but doesn't translate to a walk speedup because additive
steps re-randomise the walk in $(\ZZ/n)$.

\section{The Howe-gluing structural fact}\label{sec:howe}

\begin{theorem}[Howe~\cite{howe1996}]\label{thm:howe}
For ordinary $E_1, E_2/\FF_q$ with
\begin{enumerate}[label=\textbf{(H\arabic*)}]
\item $\Hom_{\FF_q}(E_1, E_2) = 0$,
\item $E_1[2] \cong E_2[2]$ as $\FF_q$-Galois modules with Weil
      pairing,
\item $\gcd(\#E_1, \#E_2) = 1$,
\end{enumerate}
the $(2,2)$-isogenous abelian surface $(E_1 \times E_2)/\Gamma_\alpha$,
where $\Gamma_\alpha$ is the graph of any
$\FF_q$-Galois-equivariant isomorphism
$\alpha\colon E_1[2] \to E_2[2]$, is the Jacobian of a smooth
genus-2 curve $C/\FF_q$.
\end{theorem}

For $(E_1, E_2) = (E, E^t)$ of a prime-order ordinary curve, all
three conditions hold by structural argument
(Section~\ref{sec:empirical}). The explicit construction of $C$
requires Mestre's reconstruction algorithm and is the subject of a
companion note; the \emph{existence} of $C$ follows from
Theorem~\ref{thm:howe} alone.

\begin{remark}
A naive attempt at the explicit construction
$C\colon y^2 = f_1(x) \cdot f_2(x)$ (with $f_1, f_2$ the 2-torsion
polynomials of $E_1 = E$, $E_2 = E^t$) \emph{fails} to produce a
$\Jac(C)$ that is $\FF_p$-isogenous to $E \times E^t$~--- the
construction yields a genus-2 curve whose Jacobian is simple over
$\FF_p$ (verified numerically on the toy $p = 1009$). The actual
Howe-glued construction requires Mestre's reconstruction from the
Igusa invariants of $(E \times E^t)/\Gamma_\alpha$.
\end{remark}

\subsection*{The $\FF_{p^3}$ field-of-definition obstruction for secp256k1}

\begin{proposition}[$\FF_{p^3}$ obstruction]\label{prop:fp3obs}
Let $E_0\colon y^2 = x^3 + 7$ be secp256k1 over $\FF_p$
\textup{(}$p = 2^{256} - 2^{32} - 977$, $p \equiv 1 \pmod 6$\textup{)}.
Let $E_k$ be any of the three sextic twists of $E_0$ satisfying the
Howe conditions \textup{(H1)--(H3)} of Theorem~\ref{thm:howe}.  Then:
\begin{enumerate}[label=\textup{(\roman*)}]
\item\label{it:irred} $x^3 + 7$ is irreducible over $\FF_p$; hence the
      affine $2$-torsion of $E_0$ lies in
      $E_0(\FF_{p^3}) \setminus \bigl(E_0(\FF_p) \cup E_0(\FF_{p^2})\bigr)$.
\item\label{it:partner} For each glueable partner $E_k$, the polynomial
      $x^3 + b_k$ \textup{(}$b_k = 7u^k \bmod p$, $u$ a primitive sixth
      root of unity mod $p$\textup{)} is also irreducible over $\FF_p$.
\item\label{it:galois} Any Galois-equivariant isomorphism
      $\alpha\colon E_0[2] \to E_k[2]$ is not defined over $\FF_p$
      or $\FF_{p^2}$.
\item\label{it:cover} The Howe isogeny $\Jac(C) \to E_0 \times E_k$
      is defined over $\FF_{p^3}$, not over $\FF_p$ or $\FF_{p^2}$.
\end{enumerate}
\end{proposition}

\begin{proof}
From the Frobenius trace
\[
  t = p + 1 - n = 432420386565659656852420866390673177327 \approx 2^{129}
\]
and the Newton-power-sum formula
$|E_0(\FF_{p^k})| = p^k + 1 - (\alpha^k + \bar{\alpha}^k)$
($\alpha\bar{\alpha} = p$, $\alpha + \bar{\alpha} = t$), one computes:
\begin{equation}\label{eq:torsion-counts}
  |E_0(\FF_p)| \bmod 2 = 1, \quad
  |E_0(\FF_{p^2})| \bmod 2 = 1, \quad
  |E_0(\FF_{p^3})| \bmod 4 = 0.
\end{equation}
The pattern in~\eqref{eq:torsion-counts} shows no $2$-torsion over $\FF_p$ or
$\FF_{p^2}$, but full $E_0[2] \cong (\ZZ/2)^2$ at $\FF_{p^3}$.  Since
$E_0\colon y^2 = x^3 + 7$, the affine $2$-torsion $x$-coordinates are
roots of $x^3 + 7$; they lie in $\FF_{p^3}$ precisely when $x^3 + 7$
is irreducible over $\FF_p$, confirmed by
\texttt{factormod($x^3 + 7$,\,p)} returning a single degree-$3$ factor.
Item~\ref{it:partner} is the same computation applied to $x^3 + b_k$
for each glueable partner (verified for $k \in \{1, 3, 4\}$ in the
sextic-twist labeling of \S\ref{sec:empirical}).
Items~\ref{it:galois}--\ref{it:cover} follow: a Galois-equivariant
$\alpha\colon E_0[2] \to E_k[2]$ requires both $2$-torsion groups to
be simultaneously rational, which first occurs over $\FF_{p^3}$; by
Howe's theorem~\cite{howe1996} the resulting cover $\Jac(C) \to
E_0 \times E_k$ is then defined over $\FF_{p^3}$.
PARI/GP script: \texttt{secp256k1\_cm\_audit/fp3\_obstruction\_secp256k1.gp}.
\end{proof}

\begin{corollary}[Cover-attack cost for secp256k1]\label{cor:fp3cost}
Any Howe-cover attack on secp256k1 must solve DLP in
$\Jac(C)(\FF_{p^3})$.  By Gaudry's bound~\cite{gaudry2000}, this costs
\[
  O\!\left(\sqrt{|\Jac(C)(\FF_{p^3})|}\right)
  = O\!\left(\sqrt{p^3}\right)
  = O\!\left(p^{3/2}\right)
  \approx 2^{384}
\]
group operations, versus plain Pollard~$\rho$ at
$O(\sqrt{n}) \approx 2^{128}$.
The cover-based attack is $2^{256}$ times \emph{slower} than direct ECDLP.
\end{corollary}

\paragraph{Toy verification ($p = 43$).}
For the analogous toy case $p = 43$, the Howe-glued curve
$C_{43}\colon y^2 = (x^3 + 7)(x^3 + 13)$ (the two glueable-partner
twist polynomials over $\FF_{43}$) has Igusa invariants
\[
  (J_2 : J_4 : J_6 : J_8 : J_{10}) = (39 : 27 : 7 : 15 : 36)
  \quad \in \mathbb{P}(2,4,6,8,10)/\FF_{43},
\]
confirming that the weighted-projective isomorphism class is
well-defined over $\FF_{43}$.  (The $J_2$ entry carries the
conventional factor of $2$ from PARI's \texttt{igusa\_quadruple}
normalisation relative to Cardona--Quer; $J_{10} = 36$ is
normalization-independent and matches both conventions.)
Script: \texttt{secp256k1\_cm\_audit/igusa\_f43\_howe.gp}.

\section{Empirical verification on deployed curves}\label{sec:empirical}

The companion script \texttt{multi\_curve\_audit.gp} runs the full
audit on four deployed curves; results in Table~\ref{tab:curves}.

\begin{table}[h]
\centering
\small
\begin{tabular}{lcccc}
\toprule
Property & secp256k1 & P-256 & brainpoolP256r1 & Curve25519 \\
\midrule
$j(E)$ & $0$ & generic & generic & generic \\
Small-CM disc & $D = -3$ & none & none & none \\
$\Aut(E)$ & $\ZZ/6\ZZ$ & $\{\pm 1\}$ & $\{\pm 1\}$ & $\{\pm 1\}$ \\
$\rho$ speedup & $\sqrt{6}$ & $\sqrt{2}$ & $\sqrt{2}$ & $\sqrt{2}$ \\
$n$ prime? & yes & yes & yes & no $(= 8\ell)$ \\
$3 \mid (n-1)$ & yes & yes & yes & yes (for $\ell$) \\
$p$ cube mod $n$? & yes & yes & \textbf{no} & yes (mod $\ell$) \\
Howe (H1) & \checkmark & \checkmark & \checkmark & \checkmark \\
Howe (H2) & \checkmark & \checkmark & \checkmark & mixed \\
Howe (H3) & 1 & 1 & 1 & 4 \\
(2,2) cover exists? & yes & yes & yes & via $\ell$-restriction \\
Cover-based attack & no (B5) & no (B5) & no (B5) & no (B5) \\
\textbf{Verdict} & PASS & PASS & PASS & PASS \\
\bottomrule
\end{tabular}
\caption{Cross-curve structural-audit results.}\label{tab:curves}
\end{table}

\paragraph{Surprising auxiliary finding.} Among the four curves,
three (secp256k1, P-256, Curve25519) satisfy
$p^{(n-1)/3} \equiv 1 \pmod n$~--- i.e., $p$ is a cube in
$(\ZZ/n)^*$, giving embedding degree $(n-1)/3$ rather than $n-1$.
brainpoolP256r1 is the outlier. This is not security-relevant
(embedding degree remains $\geq 253$ bits in every case) but is a
notable structural coincidence among the deployed curves.

\paragraph{Curve25519's cofactor changes Howe but not the verdict.}
Curve25519 has cofactor 8, so $\gcd(n, n_{\text{twist}}) \geq 4 > 1$
and Howe (H3) fails on the full $\#E$. Restricting to the prime
subgroup of size $\ell = 2^{252} + 27742\ldots$, the framework
applies and the verdict is the same.

\section{The structural-completeness theorem}\label{sec:theorem}

\begin{theorem}[Structural Completeness for Prime-Field Ordinary ECC]
\label{thm:main}
Let $E$ be an ordinary elliptic curve over a prime field $\FF_p$,
with prime-order subgroup of size $r$ (i.e., $\#E(\FF_p) = h \cdot r$
with $h \in \{1, 2, 4, 8\}$ and $r$ prime). Let $\mathcal{A}$ be any
algorithm in the attack family
\[
   \mathcal{A} \in \{\text{$\FF_p$-isogeny transport},\;
                    \Cl(\End)\text{-orbit walking},\;
                    (N,N)\text{-cover},
\]\[
                    \text{higher-genus cover},\;
                    \text{supersingular-reduction},\;
                    \text{VFCG-}\rho\}
\]
producing a candidate value $\tilde d$ for the secret $d$. Then
\[
   T(\mathcal{A}) := \text{expected time of $\mathcal{A}$ to produce
                          correct $\tilde d$}
   \geq \sqrt{r / |\Aut(E)|}
\]
with $|\Aut(E)| \in \{2, 4, 6\}$ the size of the curve's
automorphism group.
\end{theorem}

\begin{proof}
Case analysis over $\mathcal{A}$:
\begin{itemize}
\item \textbf{$\FF_p$-isogeny transport}: by B1 (Tate), the
      transported instance has the same $r$; the algorithm reduces
      to plain $\rho$ on the transported curve, also $\sqrt{r}$ cost.
\item \textbf{$\Cl(\End)$-orbit walking}: by B2 ($h = 1$ at the
      crater for small-CM curves) or by orbit-walk-cost $\geq$ saving
      for $h > 1$, the algorithm reduces to plain $\rho$.
\item \textbf{$(N, N)$-cover via Howe-gluing}: by B4--B5, the
      cover's Jacobian has genus 2 and DLP cost
      $O(p) > O(\sqrt{r})$. For secp256k1 specifically,
      Corollary~\ref{cor:fp3cost} sharpens this to $O(p^{3/2})
      \approx 2^{384}$ because Proposition~\ref{prop:fp3obs}
      forces the cover's field of definition to be $\FF_{p^3}$.
      Algorithm cost dominated by $\Jac(C)$ DLP.
\item \textbf{Higher-genus cover via Weil restriction}: by B5--B6.
      Genus $g \geq 2$ covers have DLP cost
      $O(p^{2 - 2/g}) \geq O(p) > O(\sqrt{r})$ over prime fields;
      Diem 2011 sub-exp is structurally blocked.
\item \textbf{Supersingular-reduction attack}: by B7. The
      supersingular isogeny problem is disjoint from $\FF_p$ ECDLP;
      canonical-lift requires anomalous curves.
\item \textbf{VFCG-$\rho$}: by Section~\ref{sec:novel}. The
      $c_\gamma$-in-small-subgroup property is satisfied uniformly
      but does not yield a walk speedup.
\end{itemize}
In each case, $T(\mathcal{A}) \geq \sqrt{r}/\sqrt{|\Aut(E)|}$, which
is the plain Pollard~$\rho$ cost.
\end{proof}

\subsection*{What the theorem says and does not say}

\textbf{Says}:
\begin{itemize}
\item For every deployed ECC curve in the prime-field ordinary
      regime, the seven specific isogeny-graph attack families do
      not yield speedups over plain $\rho$.
\item The result holds under standard assumptions: Tate's theorem,
      Howe's theorem, Gaudry's index-calculus complexity,
      Diem-Thomé's refinement, and the standing assumption that
      prime-field hyperelliptic DLP is not sub-exponentially solvable.
\end{itemize}

\textbf{Does not say}:
\begin{itemize}
\item That ECDLP is provably hard (no such result is known under
      any natural complexity assumption).
\item That side-channel / HNP-style attacks (scalar-side) are
      inapplicable.
\item That post-quantum quantum-Shor attacks fail.
\item That no future cryptanalytic technique will fall outside the
      seven enumerated families.
\end{itemize}

\section{Open questions}

\subsection*{In-scope: refinements of the theorem}
\begin{itemize}
\item Tightening B5: is the $O(p^{2 - 2/g})$ Gaudry bound tight?
\item Beyond hyperelliptic: do non-hyperelliptic genus-$g$ curves
      admit faster DLP algorithms over $\FF_p$?
\item Indirect uses of supersingular reductions for fault-attack-
      amplified ECDLP recovery on isogenous siblings: speculative.
\item Exact Cardona--Quer normalisation of Igusa invariants for
      $C_{43}$ (cross-verification with Sage / Magma pending).
\end{itemize}

\subsection*{Out-of-scope: separate attack programmes}
\begin{itemize}
\item Side-channel / HNP --- where every real-world ECDSA break of
      the past 15 years lives. Includes the proposed novel
      GLV-aware HNP variant.
\item Quantum cryptanalysis (Shor, Grover).
\item SQIsign / supersingular-isogeny-based PQC cryptanalysis.
\item Implementation cryptanalysis (fault, parsing, BIGNUM bugs).
\end{itemize}

\section{Reproducibility}

All scripts are PARI/GP 2.17+ scripts with committed expected
outputs. Total runtime approximately 3 minutes on an M-series Mac.

\begin{center}
\small
\begin{tabular}{lr}
\toprule
File & LOC \\
\midrule
\texttt{audit.gp} & 235 \\
\texttt{cl\_o\_horizontal.gp} & 203 \\
\texttt{structural\_invariants.gp} & 292 \\
\texttt{vfcg\_experiment.gp} & 254 \\
\texttt{vfcg\_variants.gp} & 284 \\
\texttt{howe\_gluing\_test.gp} & 260 \\
\texttt{howe\_explicit\_cover.gp} & 274 \\
\texttt{cover\_complexity.gp} & 165 \\
\texttt{supersingular\_reductions.gp} & 171 \\
\texttt{p256\_comparison.gp} & 283 \\
\texttt{multi\_curve\_audit.gp} & 274 \\
\texttt{curve\_audit\_template.gp} & 120 \\
\texttt{generate\_cm\_curve.gp} & 80 \\
\texttt{glv\_hnp\_phase1.gp} & 180 \\
\texttt{fp3\_obstruction\_secp256k1.gp} & 200 \\
\texttt{igusa\_f43\_howe.gp} & 58 \\
\midrule
\textbf{Total} & \textbf{$\approx 3533$} \\
\bottomrule
\end{tabular}
\end{center}

In addition, a Rust integration test
(\texttt{tests/curve\_audit.rs}, $\approx$190 lines) runs the audit
via \texttt{cargo test} on every \texttt{CurveParams} constant
defined in the codebase, asserting structural completeness for each.

\begin{thebibliography}{99}

\bibitem{tate1966} J.~Tate, ``Endomorphisms of abelian varieties
over finite fields'', \emph{Invent. Math.} 2 (1966), 134--144.

\bibitem{kohel1996} D.~Kohel, \emph{Endomorphism rings of elliptic
curves over finite fields}, Ph.D. thesis, U.~C.~Berkeley, 1996.

\bibitem{fouquetmorain2002} M.~Fouquet, F.~Morain, ``Isogeny
volcanoes and the SEA algorithm'', in \emph{ANTS-V}, LNCS 2369,
Springer (2002), 276--291.

\bibitem{howe1996} E.~W.~Howe, ``Constructing distinct curves with
isomorphic Jacobians in characteristic zero'', \emph{Israel J. Math.}
93 (1996), 1--16.

\bibitem{gaudry2000} P.~Gaudry, ``An algorithm for solving the
discrete log problem on hyperelliptic curves'', in \emph{Eurocrypt
2000}, LNCS 1807, Springer (2000), 19--34.

\bibitem{diem2011} C.~Diem, ``On the discrete logarithm problem in
elliptic curves over non-prime finite fields and on hyperelliptic
curves'', \emph{Acta Arith.} 147 (2011), 75--104.

\bibitem{smart1999} N.~P.~Smart, ``The discrete logarithm problem
on elliptic curves of trace one'', \emph{J. Cryptol.} 12 (1999),
193--196.

\bibitem{mov1993} A.~Menezes, T.~Okamoto, S.~Vanstone, ``Reducing
elliptic curve logarithms to logarithms in a finite field'',
\emph{IEEE Trans. Inf. Theory} 39 (1993), 1639--1646.

\bibitem{castryckdecru2023} W.~Castryck, T.~Decru, ``An efficient
key recovery attack on SIDH'', in \emph{Eurocrypt 2023}, LNCS
14008, Springer (2023), 423--447.

\bibitem{couveignes1997} J.-M.~Couveignes, ``Hard homogeneous
spaces'', preprint, 1997.

\bibitem{rostovtsevstolbunov2006} A.~Rostovtsev, A.~Stolbunov,
``Public-key cryptosystem based on isogenies'', IACR ePrint
2006/145.

\bibitem{csidh2018} W.~Castryck, T.~Lange, C.~Martindale, L.~Panny,
J.~Renes, ``CSIDH: an efficient post-quantum commutative group
action'', in \emph{AsiaCrypt 2018}, LNCS 11274, Springer (2018),
395--427.

\bibitem{mestre1991} J.-F.~Mestre, ``Construction de courbes de
genre 2 à partir de leurs modules'', in \emph{Effective Methods in
Algebraic Geometry}, Birkhäuser (1991), 313--334.

\bibitem{boneh1996} D.~Boneh, R.~Venkatesan, ``Hardness of computing
the most significant bits of secret keys in Diffie-Hellman and
related schemes'', in \emph{Crypto 1996}, LNCS 1109, Springer
(1996), 129--142.

\bibitem{gallant2001} R.~Gallant, R.~Lambert, S.~Vanstone, ``Faster
point multiplication on elliptic curves with efficient
endomorphisms'', in \emph{Crypto 2001}, LNCS 2139, Springer (2001),
190--200.

\end{thebibliography}

\end{document}
